% META: {"id": "1SPE-TRIGO-EV-A", "chapitre": "1SPE-TRIGONOMETRIE", "type_objet": "evaluation", "version": "A", "duree_min": 55, "bareme_total": 20, "capacites_codes": ["C1","C2","C3","C4","C5"], "competences": ["calculer","raisonner","representer","communiquer"], "mode_creation": "ex_nihilo", "sources_inspiration": [], "parametres_sympy": {}, "fichier_tex": "chapitres/1SPE-TRIGONOMETRIE/evaluations/1SPE-TRIGO-EV-A.tex", "corrige_tex": "chapitres/1SPE-TRIGONOMETRIE/evaluations/1SPE-TRIGO-EV-A-corrige.tex", "status": "generated"}
% BEGIN-VERIFY
% from sympy import *
% x = symbols('x', real=True)
% # Ex1: conversions
% assert 225 * pi / 180 == 5*pi/4
% assert Rational(11,4) - 2 == Rational(3,4)
% # Ex2: valeurs remarquables
% assert cos(5*pi/6) == -sqrt(3)/2
% assert sin(7*pi/4) == -sqrt(2)/2
% # Ex3: formule addition
% assert simplify(cos(5*pi/12) - (sqrt(6)-sqrt(2))/4) == 0
% # Ex4: equation
% assert cos(pi/3) == Rational(1,2)
% # Ex5: fonction
% f = 2*cos(x) + 1
% assert f.subs(x, 0) == 3
% assert f.subs(x, pi) == -1
% assert f.subs(x, 2*pi/3) == 0
% assert f.subs(x, 4*pi/3) == 0
% END-VERIFY

%---------------------------------------------------------------
% EVALUATION A (55 min) — Trigonometrie
% Bareme total : 20 points
%---------------------------------------------------------------

\section*{Evaluation — Trigonometrie (Version A)}
\textit{Duree : 55 min. Calculatrice interdite. Toutes les reponses doivent etre justifiees.}

\bigskip

%===============================================================
\subsection*{Exercice 1 \hfill (4 points) — C1}
%===============================================================

\begin{enumerate}
  \item Convertir $225°$ en radians. \hfill (1 pt)
  \item Convertir $\dfrac{7\pi}{6}$ en degres. \hfill (1 pt)
  \item Determiner la mesure principale de $\dfrac{11\pi}{4}$. \hfill (1 pt)
  \item Placer sur le cercle trigonometrique le point associe a $-\dfrac{2\pi}{3}$ et donner ses coordonnees. \hfill (1 pt)
\end{enumerate}

%===============================================================
\subsection*{Exercice 2 \hfill (4 points) — C2}
%===============================================================

\begin{enumerate}
  \item Donner les valeurs exactes de $\cos\!\left(\dfrac{5\pi}{6}\right)$ et $\sin\!\left(\dfrac{5\pi}{6}\right)$. \hfill (1,5 pt)
  \item Donner la valeur exacte de $\sin\!\left(\dfrac{7\pi}{4}\right)$. \hfill (1 pt)
  \item Sachant que $\cos x = -\dfrac{3}{5}$ et $x \in \left]\dfrac{\pi}{2};\pi\right[$, calculer $\sin x$. \hfill (1,5 pt)
\end{enumerate}

%===============================================================
\subsection*{Exercice 3 \hfill (4 points) — C3}
%===============================================================

\begin{enumerate}
  \item Calculer $\cos\!\left(\dfrac{5\pi}{12}\right)$ en ecrivant $\dfrac{5\pi}{12} = \dfrac{\pi}{4} + \dfrac{\pi}{6}$. \hfill (2 pts)
  \item Montrer que $\cos(2x) = 1 - 2\sin^2 x$. \hfill (1 pt)
  \item Calculer $\sin^2\!\left(\dfrac{\pi}{8}\right)$. \hfill (1 pt)
\end{enumerate}

%===============================================================
\subsection*{Exercice 4 \hfill (4 points) — C4}
%===============================================================

\begin{enumerate}
  \item Resoudre dans $\mathbb{R}$ : $\cos x = \dfrac{1}{2}$. \hfill (1,5 pt)
  \item Resoudre sur $[0;2\pi]$ : $\sin x = -\dfrac{\sqrt{2}}{2}$. \hfill (1,5 pt)
  \item Resoudre sur $[0;2\pi]$ : $\cos x \leqslant -\dfrac{1}{2}$. \hfill (1 pt)
\end{enumerate}

%===============================================================
\subsection*{Exercice 5 \hfill (4 points) — C5}
%===============================================================

On pose $f(x) = 2\cos x + 1$.

\begin{enumerate}
  \item Quelle est la periode de $f$ ? La fonction est-elle paire ? \hfill (1 pt)
  \item Dresser le tableau de variations de $f$ sur $[0;2\pi]$. \hfill (1,5 pt)
  \item Resoudre $f(x) = 0$ sur $[0;2\pi]$. \hfill (1 pt)
  \item En deduire le signe de $f(x)$ sur $[0;2\pi]$. \hfill (0,5 pt)
\end{enumerate}
